\documentclass[10pt]{article}
\usepackage[utf8]{inputenc}
\usepackage[T1]{fontenc}
\usepackage{amsmath}
\usepackage{amsfonts}
\usepackage{amssymb}
\usepackage[version=4]{mhchem}
\usepackage{stmaryrd}
\usepackage{graphicx}
\usepackage[export]{adjustbox}
\graphicspath{ {./images/} }

\begin{document}
к-рая накрывает точку $p$; соответствие $p \rightarrow p^{\prime}$ биективно). Геометрия $V(p)$ переносится на $V^{\prime \prime}\left(p^{\prime}\right)$ тождественно, т. е. расстояние между точками $p_{1}^{\prime}, p_{2}^{\prime} \in V^{\prime}$ принимается равным расстоянию между их прообразами. Вместе с тем определяется глобальная геометрия универсальной накрывающей, но она, как правило, будет совершенно отличаться от глобальной геометрии накрываемой.

На псевдосфере $U$ меридианы, уходя в бесконечность, сближаются; кроме того, они ортогональны краю $U$. Поэтому край универсальной накрывающей есть кривая в плоскости Лобачевского, к-рая является ортогональной траекторией пучка параллельных по Лобачевскому (сближающихся) прямых. Эта кривая в геометрии Лобачевского наз. о р и ц и к л о м. Таким образом, с точки зрения своей внутренней геометрии универсальная накрывающая псевдосферы $U$ есть внутренняя область орицикла в плоскости Лобачевского. Универсальная накрывающая $U^{\prime}$ псевдосферы $U$ удобнее для построения модели геометрии Лобачевского, чем сама псевдосфера $U$. Напр., сколь угодно большой круг можно расположить на $U^{\prime}$, но в $U$ слипком больпой круг без перекрытий может не уместиться. Но даже $U^{\prime}$ недостаточно обширна, чтобы вместить любой объект геометрии Лобачевского, напр. полная прямая в $U^{\prime}$ не умещается. Иначе говоря, построить с помощью $U^{\prime}$ модель всей плоскости Лобачевского нельзя.

Д. Гильбертом (1900) был поставлен вопрос: не суцествует ли поверхности, внутренняя геометрия к-рой в целом совпадает с геометрией плоскости Лобачевскоro? Из простых соображений следует, что если такая поверхность есть, то она должна иметь постоянную отрицательную кривизну и быть полной. Далее, она должна быть односвязной, поскольку односвязна плоскость Лобачевского. Но на самом деле достаточно было бы найти любую полную поверхность постоянной отрицательной кривизны (согласно изложенному выше, универсальная накрывающая такой поверхности полна и односвязна). При этом имеются в виду только регулярные поверхности (хотя бы дважды непрерывно диф. ференцируемые в любой точке, т. е. класса $C^{2}$ ).

Уже в 1901 Д. Гильберт решил свою проблему (см. [2]), причем в отрицательном смысле: в трехмерном евклидовом пространстве не существует полной регулярной поверхности постоянной отрицательной кривизны. Эта теорема привлекала внимание геометров на протяжении ряда десятилетий и привлекает внимание до сих пор. Дело в том, что с ней и с ее доказательством связано много интересных вопросов (см. ниже). При этом нали чие особых линий у поверхностей вращения постоянной отрицательной кривизны (см. рис. 4,6 ) не случайно. Оно соответствует теореме Гильберта.

Доказательство теоремы Гильберта о поверхностях любого топологич. строения прежде всего интересно тем, что в нем существенно использовано понятие универсальной накрывающей поверхности. Оно сводит дело к поверхности с простейшей топологией, именно, к односвязной поверхности. По-видимому, теорема Гильберта была одним из ранних предложений математики, где это понятие уяснилось. После публикации первой статьи Д. Гильберта с доказательством его теоремы выражалось сомнение в справедливости этого доказательства. Замечания относились к необозримости поведения асимптотич. линий, к-рые по ходу дела рассматривались на поверхности ввиду сложности ее топологии. Предполагалось даже, что Д. Гильберт перерабатывал свое доказательство, чтобы уточнить его. Между тем доказательство Д. Гильберта с самого начала было безупречным. Вероятно, этих замечаний не было бы, если бы Д. Гильберт отчетливо сказал, что исследуется полное односвязное многообразие постоянной отрицательной кривизны, т. е. сама плоскость Лобачевского, в предположении, что она изометрично и регулярно погружена в трехмерное евклидово пространство.

Интересным является также то, что в доказательстве Д. Гильберта неожиданно появляются q е б ы ш вск и е с е т и: так называется сеть, в каждом сетевом четырехугольнике к-рой противоположные стороны имеют равные длины (см. рис. 7). И вот оказывается, что асимптотич. сеть (к-рая определена и невырождена во всех точках на любой О.к. п.) в случае постоянной О. к. п.- чебышевская. Эту сеть можно принять в качестве координатной: $u=$ const, $v=$ const. Тогда линейный элемент получаст

\begin{center}
\includegraphics[max width=\textwidth]{2023_07_08_2391333da740af6080a5g-1}
\end{center}

Рис. 7 . чебышевский вид

\[
d s^{2}=d u^{2}+2 \cos \omega d u d v+d v^{2},
\]

где $\omega$ - сетевой угол. Предполагая; что поверхность регулярна, полна и односвязна, и, пользуясь чебышевским характером сети, можно доказать, что $u, v$ принимают все значения: $-\infty<u<\infty,-\infty<v<\infty$ и әтим исчерпывают сеть; иначе говоря, асимптотич. сеть в целом гомеоморфна декартовой сети на евклидовой координатной плоскости (см. [9]). С другої стороны, в предположении, что гауссова кривизна постоянна, напр. $K=-1$, из (1) получается уравнение

\[
\frac{\partial^{2} \omega}{\partial u \partial v}=\sin \omega
\]

при этом из геометрич. соображений следуют ограничения

\[
0<\omega<\pi
\]

Д. Гильберт показал, что уравнение (2) не может иметь регулярного на всей координатной плоскости репения $\omega=\omega(u, v)$, подчиненного условию (3), и тем самым полная плоскость Лобачевского не может быть регулярно и изометрично погружена в трехмерное евклидово пространство. Отсюда, как объяснялось выше, вытекает более общая формулировка этой теоремы: в трехмерном евклидовом пространстве невозможна полная регулярная постоянной О. к. п. (никакой топологич. структуры).

Во многих разделах математики случалось, что какое-нибудь утверждение или какой-нибудь вопрос играли особую роль в развитии всего раздела. В дифференциальной геометрии одним из таких утверждений была (и еще остается) теорема Гильберта. Ей посвящено много публикаций. Интерес к этой теореме вызван прежде всего сочетанием в ее доказательстве разнообразных идей и понятий (универсального накрытия, чебышевской сети). Позднее обнаружилась ее связь с физикой.

Оказалось, что уравнение (2), написанное в несколько ином виде и без ограничения (3), выражает условие на потенциал в т. н. эффекте Джозефсона из теории сверхпроводимости. Физики дали уравнению (2) название «синус-Гордона». Используя связь уравнения (2) с теорией О. к. П. (см. [13], [14]), находят в конечном виде нек-рые его решения, регулярные на всей плоскости.

С. Э. Кон-Фоссен высказал (см. [15]) прөдположение что имеет место гораздо более общая теорема, чем теорема Гильберта, именно, что условие на гауссову кривизну $K=$ const $<0$ можно заменить условием $K \leqslant$ $\leqslant$ const $<0$. Однако до нач. 60 -х гг. 20 в. не было найдено ни одной теоремы, к-рая включала бы теорему Гильберта как частный случай. Только в 1961 (см. [16]) была доказана корректность теоремы Гильберта, т. е.


\end{document}